\section{Weak-$C$ IH: the open testing axis}
\label{sec:weak-ih-open-axis}
This argument is reconstructed on 19 October 2026.  Its algebraic identities
are replayed by \texttt{check\_weak\_ih\_open\_axis.py}; it does not rely on an
unrecovered historical verification receipt.

\paragraph{Domain and source input.}
The source and testing real closed fields are independent.  Let $B_c,B_t$
be their rings of bounded elements and use $B=B_c\otimes_{\mathbb R}B_t$
for Gram coefficients.  Source square roots and testing square roots are
always taken in their respective fields.  In the weak-$C$ IH source chart,
write $z=u_c^2$ infinitesimal, $y=v_c^2$ unbounded,
$H=zh\in(0,1)$, $B_*=zb$ unbounded and $\rho=B_*/y$ bounded.
The source entry calculation gives
\[
 0\leq\rho_0\leq1,\qquad
 (B_*C)_0=\left(\frac{AH}{B_*}\right)_0
 =\left(\frac{\alpha(1-H)}{B_*}\right)_0
 =c:=\frac1{2(1+\rho_0)}.
\tag{OA1}
\]
In particular $C$ is infinitesimal and $A,\alpha$ are unbounded.
For $0<\rho_0<1$,
\[
 D_0=\frac{\rho_0^2}{2(1-\rho_0^2)};
\tag{OA2}
\]
for $\rho_0=0$, $D$ is positive infinitesimal; for $\rho_0=1$, $D$ is
unbounded.  These statements follow from
$D\sim\rho^2/[2e(1+\rho)]$ and $e_0=1-\rho_0$.
No lower bound on $n\geq0$ is imposed.

\paragraph{Statement.}
For every testing $u$ with nonzero finite residue and every infinitesimal
$w=1/v$, including $w=0$, the complete reciprocal-$v$ canonical C2 parent
is a tensor sum of squares for every state.  The same assertion holds for
complex states by applying the real quadratic form to real and imaginary
parts separately.

\paragraph{Exact parent and retained squares.}
Put $D_u=1+u^2$, $D_w=1+w^2$ and
\[
 M=W-wX+uY-uwZ,\qquad J=uwW+uX-wY-Z.
\]
The parent is
\begin{align*}
\widehat F={}&4nD_uD_wX^2+4mM^2+4CD_w(uW-Y)^2
 +4AD_uD_wY^2+4DJ^2+4\alpha D_u(Y-wZ)^2\\
&+2\{J^2+D_u(Y-wZ)^2+D_w(X+uZ)^2-D_uD_wZ^2\}.
\end{align*}
Retain the entire $n,m$ packets. When $\rho_0=0$, retain the entire
$D$ packet as well and set $d=0$. When $0<\rho_0<1$, set $d=D$.
When $\rho_0=1$, choose a sufficiently large positive real $d$ below $D$
and retain $4(D-d)J^2$.  Denote the remainder by $F_*$.  All retained
terms are tensor squares or sums of tensor squares.

Use the graph $X=0$, $Y=HwZ$, and define
\[
 P=\frac{uW}{\sqrt{B_*}},\quad E=\sqrt{B_*}\,wZ,
 \quad V=uwW,\quad
 E_Y=\sqrt{\frac{B_*}{H(1-H)}}(Y-HwZ).
\tag{OA3}
\]
Crucially, the exact root relation is
\[
 PE=VZ.\tag{OA4}
\]
On the graph, the $C$, $A,\alpha$, $d$ packets are respectively
\begin{align*}
 &4B_*CD_w(P-HE/B_*)^2,\\
 &\frac{4D_u}{B_*}\{AH^2D_w+\alpha(1-H)^2\}E^2,\\
 &4d\{V-(1+Hw^2)Z\}^2.
\end{align*}
The remaining core is exactly
\[
 2V^2-4(1+Hw^2)PE
 +\frac{2}{B_*}\{H^2+u^2(1-H)^2+H^2w^2\}E^2.
\tag{OA5}
\]
Thus every graph coefficient is bounded.  The last $E^2$ coefficient is
infinitesimal, uniformly for all $0<H<1$.

\paragraph{Restoring the $Y$ normal.}
Let $L=B_*/[H(1-H)]$ and $Y_*=HwZ$.
The first normal, divided by four, is
\begin{align*}
 L_Y:=\frac14\partial_YF_*\big|_{X=0,Y=Y_*}
={}&2CD_w(Y_*-uW)+2D_u\{AH D_w-\alpha(1-H)\}wZ\\
&-2dw\{V-(1+Hw^2)Z\}
 -w\{V-(1+Hw^2)Z\}+D_u(H-1)wZ.
\tag{OA6}
\end{align*}
Every coefficient of $L_Y/\sqrt L$ in $(P,E,V,Z)$ is infinitesimal
when $d$ is bounded, including $d=0$.  Indeed the only potentially finite contribution is
\[
 2D_u\frac{AH-\alpha(1-H)}{B_*}\sqrt{H(1-H)}E,
\]
and it has zero residue by (OA1); the additional $w^2$ contribution is
infinitesimal.  The other terms contain $C$, $w/\sqrt L$ or $1/B_*$.
The quadratic $E_Y$ coefficient has residue $4cD_{u,0}>0$, and the
$X E_Y$ coefficient has zero residue.  These statements remain valid
when $H$ or $1-H$ is infinitesimal.  No inverse of $d$, $H$ or $1-H$ enters a Gram coefficient
unaccompanied by its compensating factor.

\paragraph{The strict residue when $0<\rho_0<1$.}
Write $s=\sqrt{1+u_0^2}>1$ and $r=\rho_0$.
Using (OA4), split the core cross term as
$-4PE=-4aPE-4(1-a)VZ$, where $a$ is any real number satisfying
\[
 \frac{s}{s+r}<a<\min\left\{\frac{s}{1+r},\frac{s}{s-r}\right\}.
\tag{OA7}
\]
The interval is nonempty: both upper endpoints exceed $s/(s+r)$.
The residue Gram splits into the positive $E_Y$ pivot, a $(P,E)$ block
\[
 G_{PE}=\begin{pmatrix}2/(1+r)&-2a\\-2a&2s^2/(1+r)\end{pmatrix},
\]
and a $(V,Z,X)$ block
\[
 G_{VZX}=\begin{pmatrix}
 2+4d_0&-4d_0-2(1-a)&2u_0(1+2d_0)\\
 -4d_0-2(1-a)&4d_0&-4d_0u_0\\
 2u_0(1+2d_0)&-4d_0u_0&2s^2+4d_0u_0^2
 \end{pmatrix}.
\tag{OA8}
\]
The $(V,X)$ principal block is positive definite.  Its Schur complement
is positive exactly when
\[
 \frac{s}{s+r}<a<\frac{s}{s-r},
\]
as direct substitution of $d_0=r^2/[2(1-r^2)]$ shows.  Hence (OA7)
makes both blocks strictly positive.  This checks the entire $X$ normal,
not merely the scalar graph.

\paragraph{The endpoint $r=1$.}
For finite $d>0$, positivity of the $(V,Z,X)$ block is equivalent to
$a_-<a<a_+$, where
\[
 a_\pm(d)=
 \frac{(2d+1)s^2\pm s\sqrt{2d(2d+1)}}{(2d+1)u_0^2+1}.
\]
As $d\to\infty$, the lower endpoint tends to $s/(s+1)$ and the upper
endpoint to $s/(s-1)$.  Choose first a real $a$ strictly between
$s/(s+1)$ and $\min\{s/2,s/(s-1)\}$, then a sufficiently large finite
$d$.  The interval is nonempty because $s>1$.  Since source $D$ is
unbounded this cap is admissible, and the full residue is strict.

\paragraph{The endpoint $r=0$.}
Take $a=1$, $d=0$, having retained the whole $D$ packet, and omit
$Z$ from the root list. The residue is the direct sum of
\[
 2(P-E)^2+2u_0^2E^2,
 \qquad 2(V+u_0X)^2+2X^2,
 \qquad 2D_{u,0}E_Y^2.
\tag{OA9}
\]
It is strictly positive. The exact rewriting has bounded coefficients:
(OA5) uses only $P,E,V$, and the first $X$ normal from the core is
$uV+u(1-H)w^2Z=uV+u(1-H)wE/\sqrt{B_*}$.
Thus no unscaled $Z$ or inverse of $D$ remains.

\paragraph{Conclusion.}
In all three cases the original remainder has a bounded Gram matrix with
strictly positive real residue.  The bounded-Gram lifting lemma therefore
makes it a tensor sum of squares.  Adding the retained source-positive
packets proves the assertion.  The argument is valid for arbitrary $n\geq0$,
including $n=0$, and literal projective $w=0$; no comparison between source
and testing infinitesimals was used.
